\chapter{Phase 1 --- Point-Source Validation}
\label{chap:phase1}

\begin{milestone}
\textbf{Status:} shipped (commit \code{6b37dcb}).\\
\textbf{Acceptance:} Apollo~15 deep-sensor RMSE $\le$~\SI{1}{\kelvin}
for the discrete 3-layer model and $\le$~\SI{1.5}{\kelvin} for the
Hayne~(2017) model; Apollo~17 mean-$T$ at \SI{1}{\meter} within
$\pm\SI{2}{\kelvin}$; Chang'E-4 and ChaSTE surface-$T$ peaks within
order-of-magnitude of published values.\\
\textbf{Outcome:} all criteria satisfied --- see
\cref{tab:phase1-rmse}.
\end{milestone}

\section{Goal of Phase 1}

Phase 1 is the \emph{calibration anchor} of the entire pipeline.  It
answers a simple question: given the published regolith thermo-physical
properties, does the 1-D solver reproduce the only subsurface
temperatures we have ever measured on the Moon?  If it does not, no
amount of topographic or polar sophistication in later phases will
save the study.  The success criterion is therefore quantitative and
non-negotiable.

\section{Governing equation}

The vertically-propagating heat equation in a horizontally homogeneous
regolith column is
\begin{equation}
  \rho(z)\,c_p(T)\,\frac{\partial T}{\partial t}
  \;=\;
  \frac{\partial}{\partial z}\!\left[
    k(T,z)\,\frac{\partial T}{\partial z}
  \right],
  \qquad z \in [0, z_{\max}],\ t > 0.
  \label{eq:heat1d}
\end{equation}
The density $\rho$ and conductivity $k$ depend on depth, and $k$ and
$c_p$ additionally depend on temperature; the resulting
\emph{non-linear} PDE is solved implicitly to preserve stability at
the large time steps required for multi-lunation spin-up.

\subsection{Surface boundary condition}

The surface $z=0$ is a radiative--conductive balance
\citep{vasavada1999}:
\begin{equation}
  (1-A)\,Q_\odot(t)
  \;=\;
  \varepsilon\,\sigma\,T_s^{4}
  \;-\;
  k(T_s,0)\,\left.\frac{\partial T}{\partial z}\right|_{z=0},
  \label{eq:surface_bc}
\end{equation}
where $Q_\odot(t)$ is the \emph{instantaneous} insolation (normalised
such that noon at the equator equals $\Rsolar / d_{\mathrm{AU}}^2$),
$A$ is the bond albedo, and $\varepsilon$ the infrared emissivity.

\subsection{Bottom boundary condition: geothermal flux}

In contrast to earlier studies that imposed $\partial T / \partial z
= 0$ at depth, TSUKIMI sets a \emph{finite} basal flux consistent with
the site-specific geothermal heat flow measured by Apollo:
\begin{equation}
  -\,k(T,z_{\max})\left.\frac{\partial T}{\partial z}\right|_{z=z_{\max}}
  \;=\;
  \Qb,
  \qquad
  \Qb^{\text{A15}} = \SI{21}{\milli\watt\per\meter\squared},
  \quad
  \Qb^{\text{A17}} = \SI{15}{\milli\watt\per\meter\squared}.
  \label{eq:bottom_bc}
\end{equation}

\begin{warning}
Zero-flux bottom conditions cause a slow artificial warming of the
deep column over the \num{10}+ spin-up lunations, because any net
surface-absorbed flux has nowhere to escape.  This manifests as a
0.5--\SI{2}{\kelvin} bias in the \SI{1.5}{\meter} equilibrium
temperature --- enough to push the A15 Hayne~(2017) RMSE above the
\SI{1.5}{\kelvin} acceptance threshold.  Always use
\cref{eq:bottom_bc} with a site-appropriate $\Qb$.
\end{warning}

\subsection{Depth grid: geometric, never uniform}

The diurnal skin depth at lunar noon is $\sim$\SI{3}{\centi\meter},
while the annual skin depth is $\sim$\SI{1}{\meter}.  A uniform grid
fine enough to resolve the diurnal skin wastes 2--3 orders of
magnitude of computation in the deep regime.  TSUKIMI therefore uses
a geometric depth grid
\begin{equation}
  z_i = z_0\,\left(r^{i} - 1\right)/(r-1),
  \qquad i = 0, 1, \dots, N,
  \label{eq:geo_grid}
\end{equation}
with $z_0 = \SI{1}{\milli\meter}$, $r = 1.08$, and $N$ tuned so that
$z_N \ge \SI{5}{\meter}$.  This places 60\,\% of the grid cells within
the upper \SI{10}{\centi\meter} and the remaining 40\,\% exponentially
stretched to the bottom boundary.

\begin{warning}
Uniform grids are \emph{forbidden} in the TSUKIMI codebase
(\code{CLAUDE.md}, hard rule \#2).  A pull request that switches to a
uniform grid will fail the preflight check in
\code{lunar/grid.py::assert\_geometric}.
\end{warning}

\section{Regolith properties}

Two parameterisations are supported, selected at pipeline
instantiation time.

\subsection{Hayne 2017 smooth-variation model}

\citet{hayne2017} fit a three-parameter exponential profile to the
Diviner night-side brightness temperature:
\begin{equation}
  k(z) = \Kdeep - (\Kdeep - \Ksurf)\,\exp\!\left(-\frac{z}{\Hscale}\right),
  \label{eq:hayne_k}
\end{equation}
with $\Ksurf = \SI{7.4e-4}{\watt\per\meter\per\kelvin}$,
$\Kdeep = \SI{3.4e-3}{\watt\per\meter\per\kelvin}$ and $\Hscale$
ranging from \SI{0.05}{\meter} (equator) to \SI{0.11}{\meter} (polar).

\subsection{Discrete three-layer model}

The alternative closure \citep{mitchell2008} prescribes three piecewise
layers with distinct $(\rho, k)$ pairs:
\begin{equation}
  (\rho,k)(z) =
  \begin{cases}
    (\SI{1100}{\kilogram\per\meter\cubed},\ \SI{9.3e-4}{\watt\per\meter\per\kelvin}) & 0 \le z < \SI{0.02}{\meter} \\
    (\SI{1800}{\kilogram\per\meter\cubed},\ \SI{2.5e-3}{\watt\per\meter\per\kelvin}) & \SI{0.02}{\meter} \le z < \SI{0.60}{\meter} \\
    (\SI{2000}{\kilogram\per\meter\cubed},\ \SI{4.5e-3}{\watt\per\meter\per\kelvin}) & z \ge \SI{0.60}{\meter}.
  \end{cases}
  \label{eq:discrete_k}
\end{equation}

\subsection{Temperature dependence}

Both models inherit the same temperature-dependent heat capacity
\citep{hayne2017}:
\begin{equation}
  c_p(T) = \sum_{n=0}^{4} a_n T^n,
  \qquad
  \bm{a} = (-3.6125,\; 2.7431,\; 2.3616{\times}10^{-3},\;
             -1.2340{\times}10^{-5},\; 8.9093{\times}10^{-9}),
  \label{eq:cp_T}
\end{equation}
and a radiative-conduction term that raises the effective conductivity
near the hot surface,
\begin{equation}
  k_{\mathrm{rad}}(T) = k(z) \left[1 + \chi\left(\tfrac{T}{350}\right)^3\right],
  \qquad \chi \approx 2.7.
  \label{eq:k_rad}
\end{equation}

\section{Crank--Nicolson discretisation}

\cref{eq:heat1d} is discretised with second-order central differences
in space and a $\theta=1/2$ Crank--Nicolson scheme in time
\citep{crank1947}, yielding an implicit tridiagonal system solved at
each step by the Thomas algorithm.  The node equations read
\begin{equation}
  a_i T_{i-1}^{n+1} + b_i T_i^{n+1} + c_i T_{i+1}^{n+1}
  \;=\;
  d_i\!\left(T_{i-1}^{n}, T_i^{n}, T_{i+1}^{n}\right),
  \label{eq:cn_stencil}
\end{equation}
with coefficients updated each step from the current temperature field
because $k$ and $c_p$ depend on $T$.  The full kernel is ${\sim}120$
lines of Python accelerated with Numba's \code{@njit(cache=True)}
decorator; a representative excerpt appears in
\cref{lst:solver-core}.

\begin{lstlisting}[caption={Crank--Nicolson inner loop from \code{lunar/solver\_cn.py} (abridged).},label={lst:solver-core}]
@njit(cache=True, fastmath=True)
def step_cn(T, z, k_func, cp_func, rho, dt, Q_top, Q_bot, eps, A, sigma):
    """One Crank-Nicolson step; updates T in-place."""
    N = T.shape[0]
    dz = z[1:] - z[:-1]                       # cell widths (non-uniform)
    # Thomas coefficients
    a = np.empty(N); b = np.empty(N); c = np.empty(N); d = np.empty(N)
    for i in range(1, N - 1):
        k_m = 0.5 * (k_func(T[i - 1], z[i - 1]) + k_func(T[i], z[i]))
        k_p = 0.5 * (k_func(T[i],     z[i])     + k_func(T[i + 1], z[i + 1]))
        alpha_m = k_m * dt / (rho[i] * cp_func(T[i]) * dz[i - 1] * 0.5 * (dz[i-1]+dz[i]))
        alpha_p = k_p * dt / (rho[i] * cp_func(T[i]) * dz[i]     * 0.5 * (dz[i-1]+dz[i]))
        a[i] = -0.5 * alpha_m
        c[i] = -0.5 * alpha_p
        b[i] = 1.0 + 0.5 * (alpha_m + alpha_p)
        d[i] = (T[i]
                + 0.5 * alpha_m * (T[i - 1] - T[i])
                + 0.5 * alpha_p * (T[i + 1] - T[i]))
    # Radiative surface (Eq. 2.2): Newton-linearised on T[0]
    # ... (omitted for brevity)
    # Geothermal bottom (Eq. 2.3)
    a[-1] = -1.0; b[-1] = 1.0; c[-1] = 0.0
    d[-1] = Q_bot * dz[-1] / k_func(T[-1], z[-1])
    thomas(a, b, c, d, T)                     # overwrites T
    return T
\end{lstlisting}

\section{Spin-up protocol}

The solver is initialised with the all-column equilibrium
$T_{\mathrm{eq}} = \left[(1-A)Q_{\odot,\mathrm{mean}}/(\varepsilon\sigma)
\right]^{1/4}$ and then marched for \emph{at least} 10 synodic lunations
until the maximum per-cell change between consecutive lunations
satisfies $\max_i |\Delta T_i| < \SI{0.01}{\kelvin}$.  The synodic (not
sidereal) period \SI{29.530589}{\day} is the correct phase basis for
Apollo folding; see \code{APOLLO\_DIURNAL\_FIX\_NOTES.md} in the
\code{phase1\_validation/} folder.

\section{Apollo HFE validation}

\subsection{Data}

Apollo~15 and~17 each emplaced a pair of probes with four thermocouple
sensors per probe (plus two differential sensors, unused here).  The
tables used in this study are the PDS Geosciences archive files
\code{a15.probe.1.table.asc} and \code{a15.probe.2.table.asc} (likewise
for~17), containing ${\sim}\SI{4000}{\day}$ of cadence-decimated records.

\subsection{Stability-window detection}

Raw HFE sensors exhibit a slow secular drift that is mostly instrument
calibration, not regolith physics.  The pipeline therefore restricts
the equilibrium fit to a \emph{stability window} where the linear
slope satisfies $|dT/dt| \le \SI{0.08}{\kelvin\per\year}$ and at least
\SI{20}{\percent} of the record qualifies.  The detector is the
\code{find\_stable\_window} function of \code{apollo\_helpers.py}
(\cref{lst:stability}).

\begin{lstlisting}[caption={Stability-window detector from \code{lunar/apollo\_helpers.py}.},label={lst:stability}]
def find_stable_window(subset, slope_thresh_K_per_year=0.08, min_frac=0.20):
    """Return (idx_lo, idx_hi, T_eq, T_std) of the longest contiguous
    window whose linear trend magnitude is below slope_thresh."""
    t = subset["t_sec"].values.astype(float)
    T = subset["T_K"].values.astype(float)
    # sliding-window linear regression, pick longest window below threshold
    best = (0, len(t) - 1, T.mean(), T.std())
    for lo in range(len(t)):
        for hi in range(lo + int(min_frac * len(t)), len(t)):
            slope, _ = np.polyfit(t[lo:hi], T[lo:hi], 1)
            slope_per_year = slope * (365.25 * 86400.0)
            if abs(slope_per_year) <= slope_thresh_K_per_year:
                if (hi - lo) > (best[1] - best[0]):
                    best = (lo, hi, T[lo:hi].mean(), T[lo:hi].std())
    return best
\end{lstlisting}

\subsection{Results}

\Cref{fig:a15-stability} shows the stability-window fit for Apollo~15
probe~1 (P1) and probe~2 (P2); the analogous figure for Apollo~17 is
\cref{fig:a17-stability}.  The per-sensor breakdown confirms that
shallow sensors (\code{TC11}, \code{TC12}) exhibit a borehole
heat-short artefact that biases them ${\sim}\SI{5}{\kelvin}$ warm ---
this is why the RMSE evaluation restricts to $z\ge\SI{80}{\centi\meter}$
(see \code{SHALLOW\_SENSOR\_DIAGNOSIS.md}).

\begin{figure}[H]
\centering
\includegraphics[width=\linewidth]{a15_stability_region.pdf}
\caption[A15 stability region]{Apollo~15 HFE stability-window fit.
  Coloured bands mark the contiguous intervals where
  $|dT/dt|\le\SI{0.08}{\kelvin\per\year}$; the per-sensor equilibrium
  temperatures are the averages over these bands.}
\label{fig:a15-stability}
\end{figure}

\begin{figure}[H]
\centering
\includegraphics[width=\linewidth]{a17_stability_region.pdf}
\caption[A17 stability region]{Apollo~17 HFE stability-window fit,
  analogous to \cref{fig:a15-stability}.}
\label{fig:a17-stability}
\end{figure}

\begin{figure}[H]
\centering
\includegraphics[width=\linewidth]{a15_timeseries.pdf}
\caption[A15 depth-coloured time series]{Apollo~15 depth-coloured
  temperature time series for all four thermocouple sensors on each
  probe.  The shallow-sensor warm bias ($z<\SI{80}{\centi\meter}$) is
  visible as the systematic offset of the magenta--orange family.}
\label{fig:a15-timeseries}
\end{figure}

\subsection{Mean-temperature profile}

\Cref{fig:a15-meanT,fig:a17-meanT} compare the HFE equilibrium
profile against the Hayne~(2017) and discrete 3-layer solver runs.

\begin{figure}[H]
\centering
\includegraphics[width=\linewidth]{a15_mean_T_profile.pdf}
\caption[A15 mean-$T$ profile]{Apollo~15: HFE-derived equilibrium
  profile (markers) vs.\ the Hayne~(2017) and discrete 3-layer solver
  runs (lines).  Residual panel at right; deep-sensor RMSE is annotated.}
\label{fig:a15-meanT}
\end{figure}

\begin{figure}[H]
\centering
\includegraphics[width=\linewidth]{a17_mean_T_profile.pdf}
\caption[A17 mean-$T$ profile]{Apollo~17 equivalent of
  \cref{fig:a15-meanT}.}
\label{fig:a17-meanT}
\end{figure}

\subsection{Diurnal cycle}

The synodic-folded diurnal surface temperature at Apollo~15
(\cref{fig:a15-diurnal}) and the per-sensor diurnal grid at
Apollo~17 (\cref{fig:a17-diurnal-grid}) reveal that the discrete
3-layer model tracks both the amplitude and the phase lag, while the
Hayne~(2017) model under-predicts the peak by
${\sim}\SI{5}{\kelvin}$ at the surface because the surface $\Ksurf$ is
tuned to Diviner night-side data, not Apollo noon.

\begin{figure}[H]
\centering
\includegraphics[width=\linewidth]{a15_diurnal_cycles.pdf}
\caption[A15 diurnal cycles]{Apollo~15 surface diurnal cycle folded
  on the synodic period.  Black markers are the binned median of the
  raw HFE surface sensor; coloured lines are the two model
  predictions.}
\label{fig:a15-diurnal}
\end{figure}

\begin{figure}[H]
\centering
\includegraphics[width=\linewidth]{a17_diurnal_sensor_grid_lst.pdf}
\caption[A17 per-sensor diurnal grid]{Apollo~17 per-sensor diurnal
  grid in local solar time.  Each panel shows one thermocouple; the
  damping and phase lag with depth are recovered by the discrete
  3-layer model.}
\label{fig:a17-diurnal-grid}
\end{figure}

\section{Chang'E-4 and ChaSTE validation}

The two point-source data sets beyond Apollo are folded in at the end
of the Phase-1 notebook:

\begin{figure}[H]
\centering
\includegraphics[width=0.85\linewidth]{change4_validation.pdf}
\caption[Chang'E-4 validation]{Chang'E-4 surface temperature at
  \SI{45}{\degree}S compared with the solver driven by the
  sinusoidal insolation proxy.  Peak surface temperature agrees with
  \citet{huang2022} within order of magnitude; the residual bias is
  attributed to the sinusoidal proxy ignoring the SPA basin topography.}
\label{fig:change4-val}
\end{figure}

\begin{figure}[H]
\centering
\includegraphics[width=0.85\linewidth]{chaste_validation.pdf}
\caption[ChaSTE flat-surface validation]{Chandrayaan-3 ChaSTE at
  \SI{69}{\degree}S under flat-surface geometry.  The
  ${\sim}\SI{40}{\kelvin}$ surface-temperature bias is the \emph{motivation}
  for Phase~2: applying \cref{eq:cos_incidence} with $s=\SI{6}{\degree}$
  closes the gap.}
\label{fig:chaste-val}
\end{figure}

\section{Quantitative outcome}

\Cref{tab:phase1-rmse} collects the headline metrics.  All critical
acceptance criteria are satisfied.

\begin{table}[H]
\centering
\small
\begin{tabular}{l l S[table-format=1.3] S[table-format=+1.3] l}
\toprule
Site & Model & {RMSE [K]} & {Bias [K]} & Criterion \\
\midrule
A15 (deep) & Discrete 3-layer & 0.915 & -0.12 & $\le\SI{1}{\kelvin}$ \quad \checkmark \\
A15 (deep) & Hayne (2017)     & 1.441 & +0.44 & $\le\SI{1.5}{\kelvin}$ \checkmark \\
A17 ($z{=}\SI{1}{\meter}$) & Discrete 3-layer & 0.665 & -0.038 & $\pm\SI{2}{\kelvin}$ \checkmark \\
A17 ($z{=}\SI{1}{\meter}$) & Hayne (2017)     & 2.602 & +1.71  & $\pm\SI{2}{\kelvin}$ (marginal) \\
CE-4 ($T_{\mathrm{peak}}$) & Discrete 3-layer & -- & O(1) agreement & order-unity \checkmark \\
ChaSTE ($T_{\mathrm{peak}}$) & Discrete 3-layer & -- & $-\SI{40}{\kelvin}$ (flat) & motivates Phase 2 \\
\bottomrule
\end{tabular}
\caption[Phase-1 headline RMSE]{Headline Phase-1 validation metrics.
  The Apollo~17 Hayne RMSE is the only marginal case; the discrete
  3-layer model --- which the pipeline treats as the primary closure ---
  passes every criterion.}
\label{tab:phase1-rmse}
\end{table}

\section{Bridge to Phase 2}

The ChaSTE residual identified in \cref{fig:chaste-val} is not a
physics failure.  It is a boundary-condition failure: the solver sees
a horizontal plane at \SI{69}{\degree}S, but the lander sits on a
\SI{6}{\degree} slope.  \Cref{chap:phase2} develops the topographic
machinery required to fix this without changing a single line of the
thermal kernel.
